\rok{Јануар 2024.}{I}

Други практични тест из Алгоритама и структура података

Други практични тест одржан 13.01.2024.

\zadatak{1}{Имплементација DFSVisit методе Depth-First-Search алгоритма}
Написати имплементацију \texttt{DFSVisit} методе истоименог алгоритма.

\textbf{Додатне напомене за решавање задатка:}
\begin{itemize}
	\item Написати алгоритам.
	\item У Main методи креирати произвољан граф. Обезбедити начин за тестирање и проверу нове функционалности, односно позив \texttt{DFS} методе.
	\item У template-у је метода \texttt{private void DFSVisit(LinkedList.Node node, int nodeIndex)} која је основна метода за рекурзивно извршавање DFS алгоритма.
\end{itemize}



\zadatak{2}{Проверавање да ли су два прослеђена string-а анаграми}
Анаграми представљају речи или фразе које су формиране од истих карактера.
Написати алгоритам временске комплексности \(O(n)\) који ће из string-ова који се прослеђују проверити да ли је реч о анаграмима.

\textbf{Додатне напомене за решавање задатка:}
\begin{itemize}
	\item Користити Hash табелу која је дата у шаблону.
	\item Имплементацију писати у оквиру закоментарисане методе:
	\begin{Verbatim}[fontsize=\footnotesize, numbers=left, xleftmargin=10mm]
public static bool isAnagram(string str1, string str2)
	\end{Verbatim}
	\item У алгоритму постоји разлика између великих и малих слова.
	\item Како бисте string конвертовали у низ карактера са којим се може вршити манипулација, користити методу \texttt{ToCharArray()} по принципу: \texttt{str.ToCharArray()}.
	\item Карактери се не морају експлицитно конвертовати у целобројни тип приликом убацивања у Hash табелу.
	\item Редослед карактера приликом исписивања није битан.
\end{itemize}

\textbf{Примери:}
\begin{itemize}
	\item \textbf{Улаз 1:} \texttt{"fried", "Fired"}
	\item \textbf{Излаз 1:} \texttt{false}
	\item \textbf{Улаз 2:} \texttt{"Temporary", "Potentially"}
	\item \textbf{Излаз 2:} \texttt{false}
	\item \textbf{Улаз 3:} \texttt{"listen", "silent"}
	\item \textbf{Излаз 3:} \texttt{true}
	\item \textbf{Улаз 4:} \texttt{"read", "deed"}
	\item \textbf{Излаз 4:} \texttt{false}
\end{itemize}



\zadatak{3}{Проналажење подстабала са парним бројем непарних чворова}
У оквиру стандардне структуре класе \texttt{BinarySearchTree} написати имплементацију методе која ће омогућити испис корена подстабала која имају паран број чворова са непарном вредношћу.

\textbf{Додатне напомене за решавање задатка:}
\begin{itemize}
	\item У оквиру класе \texttt{BinarySearchTree} креирати јавну методу са потписом:
	\begin{Verbatim}[fontsize=\footnotesize, numbers=left, xleftmargin=10mm]
public void findEvenOddsSubtree()
	\end{Verbatim}
	\item Из претходно креиране јавне методе треба да се позива приватна метода са потписом:
	\begin{Verbatim}[fontsize=\footnotesize, numbers=left, xleftmargin=10mm]
private int findEvenOddsSubtree(Node curr)
	\end{Verbatim}
	\item Приватна метода треба да садржи имплементацију алгоритма којим ће се пребројавати чворови са непарном вредношћу и за свако подстабло проверавати да ли је број чворова са непарном вредношћу паран. У случају да јесте, потребно је у конзоли исписати вредност корена овог подстабла.
	\item Листове не треба посматрати.
	\item Алгоритам \textbf{ОБАВЕЗНО} писати у оквиру класе \texttt{BinarySearchTree}.
\end{itemize}

\textbf{Примери:}

\begin{center}
\begin{minipage}{\textwidth}
\centering
\begin{forest}
for tree={
	circle,
	draw,
	thick,
	fill=gray!20,
	minimum size=8mm,
	inner sep=1pt,
	s sep=8mm,
	l sep=12mm,
	edge={-, thick},
	font=\small
}
[8
	[4
		[2
			[, phantom]
			[3]
		]
		[5]
	]
	[9
		[, phantom]
		[11]
	]
]
\end{forest}

\vspace{0.3cm}
\textbf{Слика 1.} Пример BST-а.
\end{minipage}
\end{center}

\begin{itemize}
	\item \textbf{Улаз 1:} BST са слике 1
	\item \textbf{Излаз 1:} 8, 4
\end{itemize}

\begin{center}
\begin{minipage}{\textwidth}
\centering
\begin{forest}
for tree={
	circle,
	draw,
	thick,
	fill=gray!20,
	minimum size=8mm,
	inner sep=1pt,
	s sep=8mm,
	l sep=12mm,
	edge={-, thick},
	font=\small
}
[10
	[4
		[2]
		[6
			[5]
			[9]
		]
	]
	[15
		[12
			[, phantom]
			[14]
		]
		[19
			[17]
			[20]
		]
	]
]
\end{forest}

\vspace{0.3cm}
\textbf{Слика 2.} Пример BST-а.
\end{minipage}
\end{center}

\begin{itemize}
	\item \textbf{Улаз 2:} BST са слике 2
	\item \textbf{Излаз 2:} 4, 15, 6
\end{itemize}



\sablon{Шаблон другог теста јануар 2024. група I}{2023/Januar/GrupaI/AISP2023_Test2_GrupaI.linq}
